%==============================================================================
%== template for LATEX poster =================================================
%==============================================================================
%
%--A0 beamer slide-------------------------------------------------------------
\documentclass[final]{beamer}
\usepackage[orientation=landscape,size=a0,
            scale=1.50         % font scale factor
           ]{beamerposter}
           
\geometry{
  hmargin=2.5cm, % little modification of margins
}
%\usepackage{harvard}
\bibliographystyle{ieeetr}  %{agsm} %{alpha}

%
\usepackage[utf8]{inputenc}
\usepackage{wrapfig}

\linespread{1.15}
%
%==The poster style============================================================
\usetheme{sharelatex}

%==Title, date and authors of the poster=======================================
\title
[none] % Conference
{ % Poster title
The potential of declarative programming languages to support user interface
programming: the case of Elm
}

\author{ % Authors
	Simon Buist\\
	\vspace{10mm}\ \\
	\texttt{simon.buist@gmail.com}
}
\institute
{The University of Bath} % General University
\date{\today}

\begin{document}
\begin{frame}[t]
%==============================================================================
\begin{multicols}{3}
%==============================================================================
%==The poster content==========================================================
%==============================================================================

\section{1. Abstract}

\textit{The research area of this project is concerned with investigating the cognitive
load a programmer experiences while programming. Two programming languages, Elm
and JavaScript will be
used to identify differences in cognitive load for a given task. User studies
and quantitative analysis will be done to determine these differences.}

\section{2. Motivation}

\begin{wrapfigure}{l}{0.05\textwidth}
  \vspace{-20pt}
    \begin{center}
      \includegraphics[width=0.05\textwidth]{elmLogo.pdf}
    \end{center}
   \vspace{-10pt}
   \caption{Elm logo}
   \vspace{-20pt}
\end{wrapfigure}

Evan Czaplicki, in his thesis of which the Elm is the product
\cite{Czaplicki2012elm}, makes claims --- ``[Elm] makes it \textit{quick
and easy} to create and combine text, images, and video into rich
multimedia displays''. While the evaluation of Elm's usability is not the
focus of the thesis, rather, it is to establish a context for Functional
Reactive Programming and describe the implementation details, he makes
other usability claims without evidence --- ``[non-declarative frameworks
for graphical user interfaces] mire programmers in the many small,
nonessential details of handling user input and modifying the display.'',
``FRP makes GUI programming much more manageable'', and in a section
entitled \textit{The Benefits of Functional GUIs}, ``In Elm, divisions between
data code, display code, and user interaction code arise fairly
naturally, helping programmers write robust GUI code''. If these claims
are true, there is all the more evidence that Elm should be a language
of choice for GUI programmers, but experiments must be done to determine
this.

\section{3. Background}

%``Beware of the Turing tar-pit in which everything is possible but nothing of
%interest is easy.'' \cite{PerlisTuringTarpit}

The catalyst to this project was a talk about the concept of
a highly reactive development environment --- where changes in the code
result in instantaneous updates to the runtime, `on-the-fly'. This was
presented in Bret Victor's ``Inventing on Principle'' \cite{Victor2012a}.
He makes several assertions about the 
`traditional' style of coding, one statement of which is that ``most of
the developer's time is spent looking at the code, blindly without an 
immediate connection to the thing they're making''. He argues that ``so
much of creation is discovery, and you can't discover anything if you
can't see what you're doing'' --- alluding to his earlier statement that
the compile--run--debug cycle is much like this. Perhaps this rapid development
cycle is not always suitable --- in their 2013 paper \cite{Lopez2012a} show that
novices tend to ``thrash'' about, trying out many ideas that may or may not be a
solution, and executing ``poorly directed, ineffective problem solving \ldots
failing to realise they are doing it in good time, and fail to break out of
it'', whereas experts think much more about the problem at hand before
proceeding with a solution.

\section{4. Experimental Design}

\vskip1ex
\begin{figure}
\centering
\includegraphics[width=1.05\columnwidth]{ElmWalk.pdf}
\caption{A platform game in the Elm IDE}
\end{figure}
\vskip2ex

In the Elm IDE picture, there is the Elm code on the left pane, and the runtime
on the right, which updates as the code is written. To determine whether Elm is
`easier' (for some operationalisation of easy) than JavaScript, I am designing
the same task in both Elm and Javacript and am conducting user participant
studies, video recording people completing the task of writing a small GUI. I am
then going to work through a series of research --- design --- research ---
design \ldots iterations in order to find out what may be causing the differences in
cognitive load that I have observed speculatively in preliminary studies.

The task has to be designed in a way that is not idiomatic of either language --
it is unfair to, say, choose a task that leverages the unique advantages of
Elm's \texttt{Signal} type. A line such as \texttt{main = lift asText
Mouse.position} will not be anywhere near as easy to achieve the same thing in
JavaScript. 

\section{5. Further Work}

In the interest of comparing research methodologies, I am curious as to the
reliability of thematic analysis to highlight themes within such an activity as
programming. I would like to extend the Elm IDE to log keypresses and mouse
co--ordinates and obtain a different perspective on the same user studies to see
if the themes identified correlate with the objective measures of cognitive
load.


%==============================================================================
%==End of content==============================================================
%==============================================================================

%--References------------------------------------------------------------------

\subsection{References}

\bibliography{poster}

%\bibitem{ref1} J.~Doe, Article name, \textit{Phys. Rev. Lett.}

%\bibitem{ref2} J.~Doe, J. Smith, Other article name, \textit{Phys. Rev. Lett.}

%\bibitem{web} \url{http://www.google.pl}

%--End of references-----------------------------------------------------------

\end{multicols}

%==============================================================================
\end{frame}
\end{document}
